\chapter[Widmo operatorów eliptycznych]{Widmo operatorów eliptycznych: zasada min--max, asymptotyka Weyla i oszacowania wartości własnych}
\label{chap:pde18}

\section{Cel i zakres rozdziału}
\label{sec:pde18-intro}
W rozdziale~\ref{chap:pde12} skonstruowaliśmy laplasjan Dirichleta
jako operator związany z domkniętą formą i uzyskaliśmy bazę
funkcji własnych. Teraz chcemy odpowiedzieć na pytania ilościowe:
jak duże są kolejne wartości własne, jak zależą od geometrii
obszaru i współczynników oraz co dokładnie oznacza asymptotyka
liczby modów o energii poniżej ustalonego progu? Podamy dowody
zasady min--max, oszacowania Li--Yau, prawa Weyla dla laplasjanu
Dirichleta i wzoru na wiodący człon śladu półgrupy cieplnej.

W całym rozdziale $d\ge1$, a $\Omega\subset\R^d$ oznacza
\emph{ograniczony obszar}. Zwartość zanurzenia
$H_0^1(\Omega)\hookrightarrow L^2(\Omega)$ wynika już z
ograniczoności obszaru i nie wymaga regularnego brzegu.
Istotnie, wybieramy otwartą kostkę $Q$ zawierającą
$\overline\Omega$. Przedłużenie przez zero funkcji z
$C_c^\infty(\Omega)$ zachowuje normę $H^1$, a przez
gęstość wyznacza izometryczne zanurzenie
$H_0^1(\Omega)\to H_0^1(Q)$. Z twierdzenia Rellicha--Kondrachova
~\ref{thm:pde10-Rellich-p} na kostce i ciągłości ograniczenia
$L^2(Q)\to L^2(\Omega)$ wynika żądana zwartość.
Ten argument działa także dla ograniczonych, niekoniecznie
spójnych zbiorów otwartych.

Formy symetryczne rozpatrujemy w rzeczywistych przestrzeniach
Hilberta. Widmo operatora rozumiemy jako widmo jego
zespolenia; dla rozpatrywanych operatorów samosprzężonych
rzeczywista baza funkcji własnych jest również bazą
po zespoleniu, z tymi samymi wartościami własnymi
i krotnościami. W rachunkach z transformatą Fouriera
korzystamy z jej zwykłej zespolonej postaci.
Silniejsze założenia geometryczne stawiamy jawnie tam,
gdzie potrzebujemy śladów lub twierdzeń o regularności.
Dla prawa Weyla wystarcza ograniczoność i warunek
$\lambda_d(\partial\Omega)=0$; w dowodzie zostanie
to wskazane osobno.

\begin{uwaga}[Od widma do ewolucji — zakres zastosowania]
\label{rem:pde-v43-spectral-heat}
Rozkład spektralny dotyczy operatora samosprzężonego,
niekoniecznie formalnego wyrażenia $-\operatorname{div}(A\nabla)+V$
na dowolnie przyjętej dziedzinie. W rozdziale~\ref{chap:pde20}
półgrupa $e^{-tL}$ i jej rozwinięcie w funkcje własne
są używane w przypadku \emph{autonomicznego}, dodatniego
operatora Dirichleta. Dla macierzy $A(t,x)$ zależnej od
czasu, zwłaszcza niesymetrycznej, nie wolno stosować
tego samego rozwinięcia do operatora chwilowego;
konstrukcja Galerkina wykorzysta wówczas bazę
\emph{pomocniczego} laplasjanu Dirichleta.
\end{uwaga}

\section{Od formy do operatora z dyskretnym widmem}
\label{sec:pde18-form}
\begin{definicja}[Symetryczna forma eliptyczna]
Niech $\Omega\subset\R^d$ będzie niepustym ograniczonym
zbiorem otwartym. Niech $A\in L^\infty(\Omega;\R^{d\times d})$
będzie symetryczna prawie wszędzie, a $V\in L^\infty(\Omega;\R)$.
Zakładamy, że istnieją $0<a_0\le a_1<\infty$ takie, iż
$ a_0|\xi|^2\le\xi^{\mathsf T}A(x)\xi\le a_1|\xi|^2$
dla prawie każdego $x\in\Omega$ i każdego $\xi\in\R^d$.
Na $H_0^1(\Omega)$ określamy formę
\begin{equation}\label{eq:pde18-form}
 q(u,v)=\int_\Omega \bigl(A\nabla u\cdot\nabla v+Vuv\bigr)
 \,\mathrm d\lambda_d(x),\qquad q[u]=q(u,u).
\end{equation}
Przez $L_{A,V}$ oznaczamy operator w $L^2(\Omega)$
związany z tą formą (jego samosprzężoność wykażemy poniżej): $u\in D(L_{A,V})$
wtedy i tylko wtedy, gdy $u\in H_0^1(\Omega)$ oraz istnieje
$f\in L^2(\Omega)$ takie, że $q(u,v)=(f,v)_{L^2}$
dla każdego $v\in H_0^1(\Omega)$; wówczas $L_{A,V}u=f$.
Funkcja $f$ jest jednoznacznie wyznaczona, ponieważ
$C_c^\infty(\Omega)\subset H_0^1(\Omega)$ jest gęsty w
$L^2(\Omega)$.
\end{definicja}

\begin{uwaga}[Istotność dziedziny operatora]
Symbol $-\operatorname{div}(A\nabla u)+Vu$ należy na początku
rozumieć dystrybucyjnie. Przy $A\in L^\infty$
nie można bez dowodu przyjąć $D(L_{A,V})=H^2(\Omega)\cap
H_0^1(\Omega)$. Definicja przez formę jest poprawna również dla
współczynników skokowych. W tym rozdziale \emph{obszar} to
otwarty zbiór spójny, a \emph{dziedzina operatora} to $D(L_{A,V})$.
\end{uwaga}

\begin{twierdzenie}[Samosprzężoność i zwarta rezolwenta]
\label{thm:pde18-compact}
Przy założeniach poprzedniej definicji operator $L_{A,V}$
jest samosprzężony i ograniczony z dołu. Jego widmo składa się
z rzeczywistych wartości własnych, liczonych z krotnościami,
$\lambda_1\le\lambda_2\le\cdots$, przy czym
$\lambda_k\to+\infty$. Istnieje ortonormalna baza
$L^2(\Omega)$ złożona z funkcji własnych $L_{A,V}$.
\end{twierdzenie}
\begin{proof}
Oznaczmy $M=\|V_-\|_{L^\infty}$, gdzie
$V_-=\max\{-V,0\}$. Przy $c=M+1$ forma
$q_c(u,v)=q(u,v)+c(u,v)_{L^2}$ spełnia
$q_c[u]\ge a_0\|\nabla u\|_2^2+\|u\|_2^2$.
Ponadto z symetrii $A$ i ograniczenia jej wartości własnych
wynika $|A\xi|\le a_1|\xi|$, więc
\[
 |q_c(u,v)|\le a_1\|\nabla u\|_2\|\nabla v\|_2
 +(\|V\|_\infty+c)\|u\|_2\|v\|_2.
\]
Forma jest zatem ciągła i koercywna na $H_0^1(\Omega)$.
Norma $q_c[u]^{1/2}$ jest równoważna normie $H_0^1$;
w szczególności $q_c$ jest iloczynem skalarnym
czyniącym tę przestrzeń zupełną.
Dla każdego $f\in L^2$ twierdzenie Laxa--Milgrama
wyznacza jednoznacznie $u=T f\in H_0^1$ takie, że
$q_c(u,v)=(f,v)_{L^2}$ dla wszystkich $v\in H_0^1$.
Otrzymujemy ciągły operator $T:L^2\to H_0^1\hookrightarrow L^2$.
Ponieważ zanurzenie $H_0^1(\Omega)\hookrightarrow L^2(\Omega)$
jest zwarte, operator $T$ działający w $L^2$ jest zwarty.

Z symetrii formy wynika
$(Tf,g)_{L^2}=q_c(Tf,Tg)=(f,Tg)_{L^2}$;
ponadto $(Tf,f)_{L^2}=q_c[Tf]\ge0$. Gdy $Tf=0$,
to $(f,v)_{L^2}=q_c(0,v)=0$ dla wszystkich
$v\in H_0^1(\Omega)$; z gęstości wynika $f=0$.
Zatem $\ker T=\{0\}$.
Twierdzenie spektralne dla zwartych operatorów samosprzężonych
(rozdział~\ref{chap:pde12}) daje bazę ortonormalną
$(\varphi_k)$ z $T\varphi_k=\mu_k\varphi_k$,
$\mu_k>0$, $\mu_k\downarrow0$. Z definicji $T$ mamy $q(Tf,v)=(f-cTf,v)_{L^2}$,
więc $Tf\in D(L_{A,V})$ oraz $(L_{A,V}+cI)Tf=f$.
Odwrotnie, dla $u\in D(L_{A,V})$ kładziemy
$f=L_{A,V}u+cu$. Wtedy $q_c(u,v)=(f,v)_{L^2}$,
a jednoznaczność rozwiązania daje $u=Tf$.
Zatem $D(L_{A,V})=\operatorname{ran}T$ oraz
$L_{A,V}=T^{-1}-cI$.

Współrzędne w bazie $(\varphi_k)$ dają dokładny opis
dziedziny i działania operatora. Dla $u_k=(u,\varphi_k)_{L^2}$
mamy
\[
 D(L_{A,V})=
 \left\{u\in L^2:\sum_{k\ge1}\mu_k^{-2}|u_k|^2<\infty\right\},
 \qquad
 L_{A,V}u=\sum_{k\ge1}(\mu_k^{-1}-c)u_k\varphi_k.
\]
Istotnie, $u=Tf$ wtedy i tylko wtedy, gdy współrzędne
$f_k=\mu_k^{-1}u_k$ tworzą ciąg z $\ell^2$.
W szczególności dziedzina zawiera wszystkie skończone
kombinacje funkcji własnych, jest więc gęsta.
Oznaczmy $\lambda_k=\mu_k^{-1}-c$. Ponieważ $c$ jest
stałe, warunek dziedziny jest równoważny
$\sum_k\lambda_k^2|u_k|^2<\infty$.
Operator diagonalny jest symetryczny. Jeżeli $g$ należy
do dziedziny jego sprzężenia, testowanie definicji
sprzężenia na $\varphi_k$ daje współrzędne
$(L_{A,V}^*g,\varphi_k)=\lambda_k(g,\varphi_k)$.
Ich całkowalność kwadratowa wymusza
$g\in D(L_{A,V})$, co dowodzi samosprzężoności.
Identyczny argument działa po zespoleniu.

Mamy $\lambda_k\to+\infty$. Dla każdego
$z\in\C\setminus\{\lambda_k:k\ge1\}$ odległość $z$
od zbioru wartości własnych jest dodatnia, a ciągi
$(1/(\lambda_k-z))$ oraz $(\lambda_k/(\lambda_k-z))$
są ograniczone. Operator diagonalny o mnożnikach
$1/(\lambda_k-z)$ jest więc ograniczoną odwrotnością
$L_{A,V}-zI$, przyjmującą wartości w jego dziedzinie.
To dowodzi, że nie ma innych punktów widma.
Ponadto $(L_{A,V}+cI)^{-1}=T$ jest zwarty;
z tożsamości rezolwenty wynika zwartość każdej rezolwenty.
Ograniczenie z dołu wynika także bezpośrednio
z $q[u]\ge-M\|u\|_2^2$.
\end{proof}

\section{Zasada min--max i jej konsekwencje}
\label{sec:pde18-minmax}
\begin{twierdzenie}[Zasada min--max Couranta--Fischera, wersja formowa]
\label{thm:pde18-minmax}
Niech $L_{A,V}$ będzie operatorem z
twierdzenia~\ref{thm:pde18-compact}, a $q$ jego formą.
Dla każdego $k\ge1$ zachodzi
\begin{equation}\label{eq:pde18-minmax}
 \lambda_k=\min_{\substack{E\subset H_0^1(\Omega)\\\dim E=k}}
 \ \max_{u\in E\setminus\{0\}}\frac{q[u]}{\|u\|_2^2}
 =\max_{\substack{F\subset L^2(\Omega)\\\dim F=k-1}}
 \ \inf_{\substack{u\in H_0^1(\Omega)\setminus\{0\}\\u\perp F}}
 \frac{q[u]}{\|u\|_2^2}.
\end{equation}
W drugim wyrażeniu dla $k=1$ przyjmujemy $F=\{0\}$.
\end{twierdzenie}
\begin{proof}
Niech $(\varphi_j)$ będzie ortonormalną bazą funkcji
własnych, a wartości własne niech będą uporządkowane rosnąco.
Niech $c$ będzie przesunięciem użytym w dowodzie
twierdzenia~\ref{thm:pde18-compact}. Dla każdego
$u\in H_0^1(\Omega)$ z definicji operatora i symetrii formy
mamy $q_c(u,\varphi_j)=(\lambda_j+c)(u,\varphi_j)_{L^2}$.
W konsekwencji układ
$\psi_j=(\lambda_j+c)^{-1/2}\varphi_j$ jest ortonormalny
w przestrzeni Hilberta $(H_0^1(\Omega),q_c)$.
Jest też zupełny: jeśli $q_c(u,\psi_j)=0$ dla wszystkich
$j$, to $(u,\varphi_j)_{L^2}=0$ dla wszystkich $j$,
a zupełność bazy w $L^2$ daje $u=0$.
Z tożsamości Parsevala w przestrzeni formowej wynika
dla $c_j=(u,\varphi_j)_{L^2}$, że
\[
 q_c[u]=\sum_{j\ge1}|q_c(u,\psi_j)|^2
       =\sum_{j\ge1}(\lambda_j+c)|c_j|^2.
\]
Ponadto sumy częściowe $\sum_{j=1}^n c_j\varphi_j$
zbiegają do $u$ w normie formy. Odejmując
$c\|u\|_2^2=c\sum_j|c_j|^2$, otrzymujemy
$q[u]=\sum_j\lambda_j|c_j|^2$; szereg jest bezwzględnie
zbieżny, ponieważ $|\lambda_j|\le\lambda_j+c+c$.
Maksimum ilorazu Rayleigha na każdej skończenie wymiarowej
podprzestrzeni istnieje: wystarczy ograniczyć ciągłą
funkcję $q$ do zwartej sfery jednostkowej tej podprzestrzeni
w normie $L^2$.

Jeśli $E$ ma wymiar $k$, odwzorowanie
$u\mapsto((u,\varphi_1),\ldots,(u,\varphi_{k-1}))$
z $E$ do $\R^{k-1}$ ma niezerowe jądro. Istnieje więc
$0\ne u\in E$ ortogonalne do $\varphi_1,\ldots,
\varphi_{k-1}$, a wtedy $q[u]\ge\lambda_k\|u\|_2^2$.
Przy $E_k=\operatorname{span}\{\varphi_1,\ldots,\varphi_k\}$
każdy niezerowy $u$ ma iloraz Rayleigha najwyżej
$\lambda_k$. To dowodzi pierwszej równości.

Dla dowolnego $F$ o wymiarze $k-1$ istnieje niezerowy
$u\in E_k\cap F^\perp$. Stąd
$\inf_{u\perp F}q[u]/\|u\|_2^2\le\lambda_k$.
Jeśli $F=\operatorname{span}\{\varphi_1,\ldots,
\varphi_{k-1}\}$, wspomniana suma spektralna daje
odwrotną nierówność; równość osiąga $\varphi_k$.
Otrzymujemy drugą równość, a oba zewnętrzne ekstrema
są osiągane na wskazanych podprzestrzeniach.
\end{proof}

\begin{wniosek}[Porównanie współczynników i potencjałów]
\label{cor:pde18-compare}
Niech $L_1,L_2$ będą operatorami z twierdzenia
~\ref{thm:pde18-compact}, mającymi wspólną dziedzinę formową
$H_0^1(\Omega)$, a ich symetryczne formy spełniają
$q_1[u]\le q_2[u]$ dla każdego $u\in H_0^1(\Omega)$.
Wtedy $\lambda_k(L_1)\le\lambda_k(L_2)$ dla każdego $k$.
W szczególności, gdy $A$ spełnia warunki z
\ref{sec:pde18-form}, a $v_-\le V\le v_+$ prawie wszędzie,
to
\begin{equation}\label{eq:pde18-compare}
 a_0\lambda_k(-\Delta_D)+v_-
 \le\lambda_k(L_{A,V})\le
 a_1\lambda_k(-\Delta_D)+v_+.
\end{equation}
\end{wniosek}
\begin{proof}
Dla dowolnej $k$-wymiarowej podprzestrzeni $E$ maksimum
ilorazu Rayleigha $q_1$ nie przewyższa odpowiedniego maksimum
$q_2$. Bierzemy minimum po $E$ i stosujemy
\eqref{eq:pde18-minmax}. Nierówności
$a_0\|\nabla u\|_2^2+v_-\|u\|_2^2\le q[u]
\le a_1\|\nabla u\|_2^2+v_+\|u\|_2^2$
prowadzą do~\eqref{eq:pde18-compare}.
\end{proof}

\begin{wniosek}[Monotoniczność Dirichleta względem obszaru]
\label{cor:pde18-domain}
Jeżeli $\Omega_1\subset\Omega_2$ są niepustymi ograniczonymi
zbiorami otwartymi,
a $-\Delta_{D,\Omega_i}$ oznacza laplasjan Dirichleta
w $L^2(\Omega_i)$, to
$\lambda_k(-\Delta_{D,\Omega_2})\le
\lambda_k(-\Delta_{D,\Omega_1})$ dla każdego $k$.
\end{wniosek}
\begin{proof}
Dla $u\in C_c^\infty(\Omega_1)$ przedłużenie przez zero
należy do $C_c^\infty(\Omega_2)$ i zachowuje
$\|u\|_2$ oraz $\|\nabla u\|_2$. Przez przejście
do granicy w normie $H^1$ otrzymujemy izometryczne
zanurzenie $H_0^1(\Omega_1)$ w $H_0^1(\Omega_2)$
z tymi samymi własnościami.
Minimum w~\eqref{eq:pde18-minmax} dla $\Omega_2$
biegnie zatem po większej rodzinie podprzestrzeni.
\end{proof}

\begin{przyklad}[Odcinek i prostokąt]
Na przedziale $(0,L)$, gdzie $L>0$, rozwiązujemy
$-u''=\lambda u$, $u(0)=u(L)=0$.
Z tożsamości $\lambda\|u\|_2^2=\|u'\|_2^2$ wynika
$\lambda>0$ dla niezerowej funkcji własnej:
przy zerowej energii $u$ byłaby stała i znikałaby na brzegu.
Ogólne rozwiązanie dla $\lambda>0$ ma postać
$u(x)=a\cos(\sqrt\lambda x)+b\sin(\sqrt\lambda x)$.
Pierwszy warunek daje $a=0$, a drugi
$\sqrt\lambda L=n\pi$, $n\ge1$. Stąd
$\lambda_n=(\pi n/L)^2$, a unormowane funkcje własne
$\varphi_n(x)=\sqrt{2/L}\sin(n\pi x/L)$.
Normalizacja wynika z całki $\int_0^L\sin^2(n\pi x/L)
\,\mathrm dx=L/2$. Zupełność układu sinusów w $L^2(0,L)$
wynika z zupełności układu trygonometrycznego na $(-L,L)$
po nieparzystym przedłużeniu funkcji.

Na prostokącie $\Omega=(0,L_1)\times(0,L_2)$ produkty
$\varphi_m^{(L_1)}(x_1)\varphi_n^{(L_2)}(x_2)$
spełniają warunek Dirichleta i są funkcjami własnymi
o wartościach $\lambda_{m,n}=\pi^2(m^2/L_1^2+n^2/L_2^2)$,
$m,n\ge1$. Ortonormalność wynika z twierdzenia Fubiniego,
a zupełność z zupełności układów jednowymiarowych:
jeżeli funkcja $f\in L^2(\Omega)$ jest ortogonalna do
wszystkich produktów, to dla każdego $n$ funkcja
$x_1\mapsto\int_0^{L_2}f(x_1,x_2)\varphi_n^{(L_2)}(x_2)
\,\mathrm dx_2$ należy do $L^2(0,L_1)$ i znika
prawie wszędzie. Po odrzuceniu przeliczalnej sumy
zbiorów miary zero, zupełność układu w drugim przedziale
daje $f=0$ prawie wszędzie. Gdy $L_1=L_2$ i $m\ne n$,
pary $(m,n)$ i $(n,m)$ dają tę samą wartość własną
oraz ortogonalne funkcje własne. Widmo nie musi być proste.
\end{przyklad}


\begin{stwierdzenie}[Skalowanie widma przy dylatacji obszaru]
\label{prop:pde18-scale}
Niech $\Omega\subset\R^d$ będzie ograniczonym obszarem,
$a>0$, a $a\Omega=\{ax:x\in\Omega\}$.
Dla każdego $k\ge1$ zachodzi
$\lambda_k(-\Delta_{D,a\Omega})=a^{-2}\lambda_k(-\Delta_{D,\Omega})$.
W konsekwencji
$N_{D,a\Omega}(\Lambda)=N_{D,\Omega}(a^2\Lambda)$.
\end{stwierdzenie}
\begin{proof}
Dla $v\in H_0^1(\Omega)$ określamy
$Uv(x)=a^{-d/2}v(x/a)$ dla $x\in a\Omega$.
Podstawienie $x=ay$ daje
\[
 \|Uv\|_{L^2(a\Omega)}^2
 =a^{-d}\int_{a\Omega}|v(x/a)|^2\,\mathrm d\lambda_d(x)
 =\|v\|_{L^2(\Omega)}^2.
\]
Odwrotnością jest $U^{-1}w(y)=a^{d/2}w(ay)$,
więc $U$ jest izometrycznym izomorfizmem (unitarnym
po zespoleniu). Dla funkcji gładkich o zwartym nośniku
$\nabla Uv(x)=a^{-d/2-1}\nabla v(x/a)$, skąd
$\|\nabla Uv\|_{L^2(a\Omega)}^2
=a^{-2}\|\nabla v\|_{L^2(\Omega)}^2$.
Przez gęstość i zastosowanie odwrotności wnioskujemy,
że $U$ przeprowadza $H_0^1(\Omega)$ na $H_0^1(a\Omega)$
oraz że tożsamość dla gradientu obowiązuje na całej
dziedzinie formowej.
Zasada min--max zastosowana do $U(E)$ i $U^{-1}(E)$
daje pierwszą równość, a definicja funkcji zliczającej drugą.
\end{proof}

\section{Funkcja zliczająca i prawo Weyla}
\label{sec:pde18-weyl}
\begin{definicja}[Funkcja zliczająca]
Dla operatora o dyskretnym widmie ograniczonym z dołu
i wartościach własnych $\lambda_k(L)\to+\infty$ określamy
$N_L(\Lambda)=\#\{k\ge1:\lambda_k(L)\le\Lambda\}$,
licząc krotności wartości własnych. Symbol $\omega_d$
oznacza objętość jednostkowej kuli w $\R^d$.
\end{definicja}

\begin{lemat}[Liczba modów na sześcianie]
\label{lem:pde18-cube}
Niech $Q=(0,h)^d$, $h>0$. Dla laplasjanu z warunkiem
Dirichleta albo Neumanna zachodzi
\begin{equation}\label{eq:pde18-cube}
 N_{D,Q}(\Lambda)=\frac{\omega_dh^d}{(2\pi)^d}
 \Lambda^{d/2}+O_{d,h}(1+\Lambda^{(d-1)/2}),
 \qquad \Lambda\to\infty,
\end{equation}
a dla $N_{N,Q}(\Lambda)$ obowiązuje ten sam człon główny
i oszacowanie reszty tego samego rzędu.
\end{lemat}
\begin{proof}
W przypadku Dirichleta bierzemy produkty unormowanych
sinusów z poprzedniego przykładu. Dla Neumanna na $(0,h)$
bierzemy $h^{-1/2}$ oraz $\sqrt{2/h}\cos(n\pi x/h)$,
$n\ge1$. Spełniają one zerowy warunek na pochodną
na końcach, a zupełność wynika z parzystego przedłużenia
na $(-h,h)$ i zupełności układu trygonometrycznego.
Produkty w $d$ wymiarach są zupełne przez powtarzanie
argumentu Fubiniego z przykładu prostokąta. Całkowanie
przez części w każdej zmiennej sprawdza równanie formowe
dla testów gładkich (w przypadku Neumanna aż do brzegu);
przez gęstość równanie obowiązuje dla wszystkich testów
z odpowiedniej dziedziny formowej. Zatem wszystkie
wartości własne mają postać
$\pi^2h^{-2}\sum_{j=1}^d n_j^2$, przy czym
$n\in\{1,2,\ldots\}^d$ dla Dirichleta i
$n\in\{0,1,\ldots\}^d$ dla Neumanna. Każdy indeks
wielokrotny liczymy oddzielnie, także gdy odpowiadające
mu wartości własne się pokrywają.

Niech $R=h\sqrt\Lambda/\pi$ dla $\Lambda\ge0$ oraz
$K_R=\{x\in[0,\infty)^d:|x|\le R\}$.
Dla każdego $n\in\{1,2,\ldots\}^d$ o $|n|\le R$
bierzemy półotwartą kostkę $\prod_{i=1}^d[n_i-1,n_i)$.
Kostki te są rozłączne i ich suma $S_R$ ma miarę
$N_{D,Q}(\Lambda)$. Mamy
$K_{(R-\sqrt d)_+}\subset S_R\subset K_R$
z dokładnością do zbiorów miary zero: dla punktu
w mniejszym ortancie wybieramy $n_i=\lfloor x_i\rfloor+1$,
wtedy $|n|\le|x|+\sqrt d\le R$ (dla $R<\sqrt d$
mniejszy ortant ma miarę zero). Drugie zawieranie
wynika z $0\le x_i\le n_i$.
Zatem
\[
 2^{-d}\omega_d(R-\sqrt d)_+^d
 \le N_{D,Q}(\Lambda)\le 2^{-d}\omega_dR^d.
\]
Różnica skrajnych wyrażeń jest
$O_d(1+R^{d-1})$: dla $R\ge\sqrt d$ stosujemy
twierdzenie o wartości średniej do funkcji $t^d$,
a dla mniejszych $R$ różnica jest ograniczona.
Indeksy Neumanna z co najmniej jedną zerową współrzędną
mają liczbę nie większą niż $d(\lfloor R\rfloor+1)^{d-1}$,
gdyż wybieramy zerową współrzędną, a pozostałe
należą do $\{0,\ldots,\lfloor R\rfloor\}$.
Po podstawieniu $R=h\sqrt\Lambda/\pi$ otrzymujemy
oba oszacowania, także dla $d=1$.
\end{proof}

\begin{lemat}[Rozdzielanie na sześciany i ograniczenie liczby modów]
\label{lem:pde18-bracket}
Niech $Q_1,\ldots,Q_m$ będą otwartymi sześcianami
pochodzącymi z tej samej siatki, o parami rozłącznych
wnętrzach, a $P=\operatorname{int}(\bigcup_{j=1}^m
\overline{Q_j})$, gdzie $m\ge1$. Dla każdego $\Lambda\in\R$
\begin{equation}\label{eq:pde18-bracket}
 \sum_{j=1}^m N_{D,Q_j}(\Lambda)
 \le N_{D,P}(\Lambda)
 \le\sum_{j=1}^m N_{N,Q_j}(\Lambda).
\end{equation}
\end{lemat}
\begin{proof}
Przestrzeń $L^2(P)$ utożsamiamy izometrycznie z
$\bigoplus_{j=1}^m L^2(Q_j)$ za pomocą ograniczeń,
gdyż $P\setminus\bigcup_j Q_j$ leży w skończonej sumie
hiperpłaszczyzn i ma miarę zero. Przy tym utożsamieniu
mamy zawierania dziedzin formowych
\[
 \bigoplus_{j=1}^m H_0^1(Q_j)
 \subset H_0^1(P)\subset
 \bigoplus_{j=1}^m H^1(Q_j).
\]
Pierwsze wynika z przedłużenia przez zero z każdej
komórki i skończonego sumowania, drugie z ograniczenia
funkcji Sobolewa na otwarte podzbiory. Energia na każdej
z tych przestrzeni jest dana przez sumę
$\sum_j\int_{Q_j}|\nabla u_j|^2\,\mathrm d\lambda_d$,
a kwadrat normy $L^2$ przez sumę kwadratów norm.

Dla form na skrajnych przestrzeniach bazę własną tworzy
suma układów własnych poszczególnych sześcianów,
przedłużonych przez zero jako funkcje $L^2(P)$.
W przypadku Neumanna konstrukcja z twierdzenia
~\ref{thm:pde18-compact} działa na każdej przestrzeni
$H^1(Q_j)$: po dodaniu normy $L^2$ forma jest
koercywna, a zanurzenie jest zwarte z twierdzenia
~\ref{thm:pde10-Rellich-p}. Dowód zasady min--max
pozostaje ten sam dla tych form i ich skończonych sum prostych.
Jeżeli $\lambda_k^-$, $\lambda_k^+$ oznaczają uporządkowane
wartości własne sum odpowiednio Dirichleta i Neumanna,
to zawierania dziedzin dają
$\lambda_k^+\le\lambda_k(-\Delta_{D,P})\le\lambda_k^-$.
Dla dowolnego progu $\Lambda$ oznacza to
$N^-(\Lambda)\le N_{D,P}(\Lambda)\le N^+(\Lambda)$,
gdzie $N^\pm$ są sumami funkcji zliczających komórek,
co dowodzi~\eqref{eq:pde18-bracket}.
\end{proof}

\begin{twierdzenie}[Prawo Weyla dla laplasjanu Dirichleta]
\label{thm:pde18-weyl}
Niech $\Omega\subset\R^d$ będzie ograniczonym obszarem
o brzegu miary Lebesgue'a zero, tj.
$\lambda_d(\partial\Omega)=0$. Niech $-\Delta_D$
oznacza laplasjan Dirichleta w $L^2(\Omega)$.
Wówczas
\begin{equation}\label{eq:pde18-weyl}
 \lim_{\Lambda\to\infty}
 \frac{N_{-\Delta_D}(\Lambda)}{\Lambda^{d/2}}
 =\frac{\omega_d}{(2\pi)^d}|\Omega|.
\end{equation}
Równoważnie, gdy $k\to\infty$,
\begin{equation}\label{eq:pde18-lambda-weyl}
 \lambda_k(-\Delta_D)\sim
 (2\pi)^2\left(\frac{k}{\omega_d|\Omega|}\right)^{2/d}.
\end{equation}
\end{twierdzenie}
\begin{proof}
Ustalmy siatkę sześcianów o boku $h>0$ w $\R^d$.
Ze względu na ograniczoność $\Omega$ tylko skończenie
wiele komórek siatki przecina $\Omega$. Przy dostatecznie
małym $h$ istnieje przynajmniej jedna domknięta komórka
zawarta w $\Omega$ (wybieramy kulę wewnętrzną
o promieniu większym niż $\sqrt d\,h$).
Przez $P_h^-$ oznaczmy wnętrze sumy wszystkich domkniętych
komórek zawartych w $\Omega$, a przez $P_h^+$ wnętrze sumy
wszystkich domkniętych komórek, których wnętrza przecinają
$\Omega$. Mamy $P_h^-\subset\Omega\subset P_h^+$.
Istotnie, pierwszy zbiór jest zawarty w $\Omega$
z definicji. Dla każdego $x\in\Omega$ istnieje kula
$B_\rho(x)\subset\Omega$. Każdy punkt tej kuli należy do
domknięcia komórki, której wnętrze przecina $\Omega$
(wnętrze komórki jest gęste w jej domknięciu,
a $\Omega$ jest otwarty). Cała kula znajduje się więc
w sumie domkniętych komórek wyznaczających $P_h^+$,
czyli $x$ należy do wnętrza tej sumy.

Z monotoniczności Dirichleta względem obszaru wynika
$N_{D,P_h^-}(\Lambda)\le N_{D,\Omega}(\Lambda)
\le N_{D,P_h^+}(\Lambda)$.
Stosujemy nierówności~\eqref{eq:pde18-bracket} do obu
wielościanów, a następnie~\eqref{eq:pde18-cube} do ich
skończenie wielu komórek. Dla każdego ustalonego $h$
otrzymujemy po podzieleniu przez $\Lambda^{d/2}$ i
przejściu $\Lambda\to\infty$:
\[
 \frac{\omega_d}{(2\pi)^d}|P_h^-|
 \le\liminf_{\Lambda\to\infty}\frac{N_{D,\Omega}(\Lambda)}{\Lambda^{d/2}}
 \le\limsup_{\Lambda\to\infty}\frac{N_{D,\Omega}(\Lambda)}{\Lambda^{d/2}}
 \le\frac{\omega_d}{(2\pi)^d}|P_h^+|.
\]
Z dokładnością do zbioru miary zero różnica
$P_h^+\setminus P_h^-$ jest sumą wnętrz komórek,
których wnętrza przecinają $\Omega$, ale których
domknięcia nie są zawarte w $\Omega$. Każde takie
domknięcie zawiera punkt z $\Omega$ i punkt z jej
dopełnienia; łączący je odcinek w domkniętej komórce
przecina $\partial\Omega$. Stąd wszystkie punkty
tej komórki leżą w odległości co najwyżej $\sqrt d\,h$
od $\partial\Omega$. Zatem
\[
 |P_h^+\setminus P_h^-|
 \le\lambda_d(\{x:\operatorname{dist}(x,\partial\Omega)
 \le\sqrt d\,h\}).
\]
Ponieważ $\partial\Omega$ jest zwartym zbiorem miary zero,
prawa strona dąży do zera: domknięte otoczenia maleją
do $\partial\Omega$, a pierwsze z nich ma skończoną miarę,
więc stosujemy ciągłość miary z góry.
Zatem $|P_h^-|\to|\Omega|$ oraz
$|P_h^+|\to|\Omega|$, co daje~\eqref{eq:pde18-weyl}.

Aby uzyskać~\eqref{eq:pde18-lambda-weyl}, oznaczmy
$C_\Omega=\omega_d|\Omega|/(2\pi)^d>0$. Dla
$0<\eps<C_\Omega$ i dużych $\Lambda$ mamy
$(C_\Omega-\eps)\Lambda^{d/2}\le N(\Lambda)
\le(C_\Omega+\eps)\Lambda^{d/2}$.
Dla dostatecznie dużego $k$ stosujemy dolne oszacowanie
do $\Lambda_+=(k/(C_\Omega-\eps))^{2/d}$, co daje
$N(\Lambda_+)\ge k$, a zatem $\lambda_k\le\Lambda_+$.
Górne oszacowanie zastosowane do $\Lambda=\lambda_k$
daje $k\le N(\lambda_k)\le(C_\Omega+\eps)\lambda_k^{d/2}$,
czyli
\[
 \left(\frac{k}{C_\Omega+\eps}\right)^{2/d}
 \le\lambda_k\le
 \left(\frac{k}{C_\Omega-\eps}\right)^{2/d}.
\]
Dzielimy przez $(k/C_\Omega)^{2/d}$ i przechodzimy
do granic dolnej i górnej, a następnie $\eps\downarrow0$.
Odwrotnie, jeśli $\lambda_k\sim(k/C_\Omega)^{2/d}$,
to dla dużych $k$ i dowolnego $0<\eps<1$ mamy
$(1-\eps)(k/C_\Omega)^{2/d}\le\lambda_k
\le(1+\eps)(k/C_\Omega)^{2/d}$. Dla dużego $\Lambda$
kładziemy $k=N(\Lambda)$; wtedy
$\lambda_k\le\Lambda<\lambda_{k+1}$, więc
\[
 C_\Omega\left(\frac{\Lambda}{1+\eps}\right)^{d/2}-1
 <N(\Lambda)\le
 C_\Omega\left(\frac{\Lambda}{1-\eps}\right)^{d/2}.
\]
Po podzieleniu przez $\Lambda^{d/2}$ i przejściu
$\Lambda\to\infty$, a następnie $\eps\downarrow0$,
odzyskujemy~\eqref{eq:pde18-weyl}. Oba argumenty
uwzględniają krotności, także w punktach skoku funkcji
zliczającej.
\end{proof}

\begin{uwaga}[Co pozostaje trudne w ogólnej teorii Weyla?]
Powyższy dowód wykorzystuje wyłącznie geometrię siatki,
min--max i jawne widmo sześcianu. Dla operatorów z macierzą
$A(x)$ oczekiwany człon główny dla odpowiednio regularnych
współczynników to
$(2\pi)^{-d}\omega_d\Lambda^{d/2}
\int_\Omega(\det A(x))^{-1/2}\,\mathrm d\lambda_d(x)$.
Nie wyprowadzamy tu tej mocniejszej asymptotyki, ponieważ
wymaga ona dodatkowej lokalizacji w przestrzeni fazowej.
Natomiast porównanie~\eqref{eq:pde18-compare} natychmiast
daje ilościowe dwustronne ograniczenia $\lambda_k(L_{A,V})
\asymp k^{2/d}$ dla dostatecznie dużych $k$.
\end{uwaga}

\section{Nierówność Li--Yau: dolne oszacowanie sumy wartości własnych}
\label{sec:pde18-liyau}
\begin{lemat}[Zasada wypełniania najniższych częstotliwości]
\label{lem:pde18-bathtub}
Niech $F:\R^d\to[0,M]$ będzie mierzalna,
$0<\int_{\R^d}F\,\mathrm d\lambda_d=k<\infty$,
$M>0$, i niech $\rho=(k/(M\omega_d))^{1/d}$.
Wtedy
\begin{equation}\label{eq:pde18-bathtub}
 \int_{\R^d}|\xi|^2F(\xi)\,\mathrm d\lambda_d(\xi)
 \ge\frac{d}{d+2}k\rho^2.
\end{equation}
\end{lemat}
\begin{proof}
Jeżeli moment $\int_{\R^d}|\xi|^2F\,\mathrm d\lambda_d$
jest nieskończony, nierówność jest natychmiastowa.
Zakładamy dalej, że jest skończony.
Oznaczmy $F_0=M\chi_{B_\rho(0)}$.
Na $B_\rho$ mamy $(|\xi|^2-\rho^2)(F-F_0)\ge0$,
ponieważ oba czynniki są niedodatnie; poza kulą oba są
nieujemne. Całkując i korzystając z równości
$\int_{\R^d}F\,\mathrm d\lambda_d=\int_{\R^d}F_0\,\mathrm d\lambda_d=k$, dostajemy
$\int_{\R^d}|\xi|^2F(\xi)\,\mathrm d\lambda_d(\xi)\ge\int_{\R^d}|\xi|^2F_0(\xi)\,\mathrm d\lambda_d(\xi)$.
Całkowanie radialne daje
\[
 \int_{\R^d}|\xi|^2F_0(\xi)\,\mathrm d\lambda_d(\xi)
 =Md\omega_d\int_0^\rho r^{d+1}\,\mathrm dr
 =\frac{d}{d+2}k\rho^2.
\]
Równość zachodzi wtedy i tylko wtedy, gdy $F=F_0$
prawie wszędzie: nieujemna funkcja
$(|\xi|^2-\rho^2)(F-F_0)$ musi mieć całkę zero,
a pierwszy czynnik znika tylko na sferze miary zero.
\end{proof}

\begin{twierdzenie}[Nierówność Li--Yau dla laplasjanu Dirichleta]
\label{thm:pde18-liyau}
Niech $\Omega\subset\R^d$ będzie ograniczonym obszarem,
a $0<\lambda_1\le\lambda_2\le\cdots$ wartościami własnymi
laplasjanu Dirichleta. Dla każdego $k\ge1$ zachodzi
\begin{equation}\label{eq:pde18-liyau}
 \sum_{j=1}^k\lambda_j\ge
 \frac{d}{d+2}(2\pi)^2
 \left(\frac{1}{\omega_d|\Omega|}\right)^{2/d}
 k^{1+2/d}.
\end{equation}
W szczególności
$\lambda_k\ge\frac d{d+2}(2\pi)^2
(k/(\omega_d|\Omega|))^{2/d}$.
\end{twierdzenie}
\begin{proof}
Niech $\varphi_1,\ldots,\varphi_k$ będą ortonormalnymi
funkcjami własnymi w $L^2(\Omega)$ i przedłużmy je przez
zero do $\R^d$. Ponieważ należą do $H_0^1(\Omega)$,
ich przedłużenia należą do $H^1(\R^d)$
(twierdzenie~\ref{thm:pde9-zero-extension}), a ich słabe
gradienty są przedłużeniami przez zero gradientów na $\Omega$.
Ponadto przedłużone funkcje należą do $L^1(\R^d)$, gdyż
$\|\varphi_j\|_1\le|\Omega|^{1/2}\|\varphi_j\|_2$.
Używamy unitarnej transformaty Fouriera
$\widehat f(\xi)=(2\pi)^{-d/2}
\int_{\R^d}f(x)e^{-ix\cdot\xi}\,\mathrm d\lambda_d(x)$.
Definiujemy $F(\xi)=\sum_{j=1}^k|\widehat\varphi_j(\xi)|^2$.

Z twierdzenia Plancherela wynika
$\int_{\R^d}F(\xi)\,\mathrm d\lambda_d(\xi)=k$.
Dla $f\in H^1(\R^d)$ zachodzi
$\widehat{\partial_\ell f}=i\xi_\ell\widehat f$
w $L^2$: najpierw całkujemy przez części dla
$f\in C_c^\infty(\R^d)$, następnie przybliżamy w $H^1$
i stosujemy Plancherela. Identyfikację granicy
$i\xi_\ell\widehat f$ uzasadnia zbieżność w $L^2$
na każdej ograniczonej kuli w przestrzeni częstotliwości,
gdzie mnożnik $\xi_\ell$ jest ograniczony.
Testując równanie własne funkcją $\varphi_j$, mamy
$\|\nabla\varphi_j\|_2^2=\lambda_j$.
Stąd
\begin{equation}\label{eq:pde18-fft-energy}
 \int_{\R^d}|\xi|^2F(\xi)\,\mathrm d\lambda_d(\xi)
 =\sum_{j=1}^k\int_\Omega|\nabla\varphi_j|^2
 \,\mathrm d\lambda_d(x)=\sum_{j=1}^k\lambda_j.
\end{equation}
Dla każdego ustalonego $\xi$ funkcje
$\varphi_1,\ldots,\varphi_k$ są ortonormalne także
w zespolonej przestrzeni $L^2(\Omega;\C)$.
Przy konwencji $(f,g)=\int_\Omega f\overline g\,\mathrm d\lambda_d$
mamy $\widehat\varphi_j(\xi)=(\varphi_j,g_\xi)$,
gdzie $g_\xi(x)=(2\pi)^{-d/2}e^{ix\cdot\xi}$.
Nierówność Bessela daje
\[
 F(\xi)=\sum_{j=1}^k|(\varphi_j,g_\xi)|^2
 \le\|g_\xi\|_2^2=(2\pi)^{-d}|\Omega|=:M.
\]
Stosujemy lemat~\ref{lem:pde18-bathtub} z tym $M$,
$\rho^d=k/(M\omega_d)=(2\pi)^dk/(\omega_d|\Omega|)$.
Wstawiając oszacowanie do~\eqref{eq:pde18-fft-energy},
otrzymujemy~\eqref{eq:pde18-liyau}. Ponieważ
$k\lambda_k\ge\sum_{j=1}^k\lambda_j$, otrzymujemy
podany dolny szacunek pojedynczej wartości własnej.
\end{proof}

\begin{uwaga}[Dlaczego współczynnik jest naturalny?]
Prawa strona~\eqref{eq:pde18-liyau} jest dokładnie całką
średniego kosztu $|\xi|^2$ przy równomiernym wypełnianiu
najniższych częstotliwości w przestrzeni fazowej. Jest to wariacyjna interpretacja problemu
optymalizacyjnego dla funkcji $F$ o zadanej masie i ograniczeniu
punktowym; dowód nie wymaga znajomości prawa Weyla.
\end{uwaga}


\begin{wniosek}[Oszacowanie funkcji zliczającej bez prawa Weyla]
\label{cor:pde18-count-liyau}
Niech $d\ge1$, a $\Omega\subset\R^d$ będzie ograniczonym
obszarem. Dla $\Lambda>0$ funkcja zliczająca laplasjanu
Dirichleta spełnia
\[
 N_{-\Delta_D}(\Lambda)\le
 \frac{\omega_d|\Omega|}{(2\pi)^d}
 \left(\frac{d+2}{d}\right)^{d/2}\Lambda^{d/2}.
\]
Wniosek nie wymaga założenia $\lambda_d(\partial\Omega)=0$.
\end{wniosek}
\begin{proof}
Niech $k=N_{-\Delta_D}(\Lambda)$. Gdy $k=0$, nie ma
czego dowodzić. W przeciwnym razie $\lambda_k\le\Lambda$.
Na mocy nierówności~\eqref{eq:pde18-liyau} mamy
$\Lambda\ge\lambda_k\ge\frac d{d+2}(2\pi)^2
(k/(\omega_d|\Omega|))^{2/d}$. Przekształcenie tego
oszacowania daje tezę. Zwróćmy uwagę, że jest to
\emph{nierówność dla każdej} $\Lambda>0$, a nie
asymptotyka z prawa Weyla.
\end{proof}

\section{Pierwsza wartość własna, dodatniość i prostota}
\label{sec:pde18-ground}
\begin{twierdzenie}[Prostota pierwszej wartości własnej laplasjanu]
\label{thm:pde18-simple}
Niech $\Omega\subset\R^d$ będzie ograniczonym obszarem
(nie zakładamy regularności brzegu). Dla laplasjanu Dirichleta
$\lambda_1>0$ jest prosta, a odpowiadającą jej funkcję
własną można wybrać ściśle dodatnią we wnętrzu $\Omega$.
\end{twierdzenie}
\begin{proof}
Wybieramy kostkę $Q$ zawierającą $\overline\Omega$.
Przedłużając przez zero do $Q$ i stosując wniosek
~\ref{cor:pde9-Poincare-zero}, otrzymujemy
$\|w\|_2\le C_Q\|\nabla w\|_2$ dla $w\in H_0^1(\Omega)$.
Z zasady min--max wynika więc
$\lambda_1=\min_{w\in H_0^1\setminus\{0\}}
\|\nabla w\|_2^2/\|w\|_2^2\ge C_Q^{-2}>0$.
Niech $u$ będzie rzeczywistą niezerową funkcją własną
dla $\lambda_1$. Uzasadnimy, że $|u|\in H_0^1(\Omega)$
bez odwoływania się do śladu na brzegu.
Wybieramy $u_n\in C_c^\infty(\Omega)$ zbieżne do $u$
w $H^1$. Funkcje $|u_n|$ należą do $H_0^1(\Omega)$,
gdyż mają zwarty nośnik i można je przybliżyć splotami
w $H^1$. Z reguł dla obcięć (wniosek~\ref{cor:pde9-truncation})
wynika $\|\nabla|u_n|\|_2\le\|\nabla u_n\|_2$.
Ciąg jest ograniczony w $H_0^1$, a w $L^2$ zbiega
do $|u|$ na mocy nierówności
$||u_n|-|u||\le|u_n-u|$. Słabo zbieżny podciąg
w przestrzeni Hilberta $H_0^1$ musi zatem mieć
granicę $|u|$, co daje żądane przynależenie.
Ponadto $|\nabla|u||\le|\nabla u|$ prawie wszędzie,
więc iloraz Rayleigha $w=|u|$ jest nie większy od
$\lambda_1$ i z minimalności jest równy $\lambda_1$.

Dla każdego $v\in H_0^1(\Omega)$ i $s\in\R$ mamy
$\|\nabla(w+sv)\|_2^2-\lambda_1\|w+sv\|_2^2\ge0$,
z równością w $s=0$. Pochodna tego wielomianu
w zerze znika, skąd
\[
 \int_\Omega\nabla w\cdot\nabla v\,\mathrm d\lambda_d
 =\lambda_1\int_\Omega wv\,\mathrm d\lambda_d.
\]
Zatem $w$ jest również funkcją własną.

Gładkość wewnętrzną zarówno $u$, jak i $w$ uzyskamy
z lematu Weyla~\ref{thm:pde7-weyl}, bez założeń o brzegu
i bez osobnego argumentu dla $d=1$.
Jeżeli $f$ jest dowolną z tych funkcji, określamy
$H(x,t)=f(x)\cosh(\sqrt{\lambda_1}t)$ na
$\Omega\times\R$. Funkcja jest lokalnie całkowalna
i dystrybucyjnie harmoniczna:
dla $\Phi\in C_c^\infty(\Omega\times\R)$ całkujemy
najpierw w $x$, używając
$\Delta_x f=-\lambda_1 f$, a dla członu
$\partial_t^2\Phi$ całkujemy dwukrotnie przez części w $t$.
Człony $-\lambda_1H$ i $\lambda_1H$ się znoszą.
Lemat Weyla w wymiarze $d+1\ge2$ daje gładkiego
reprezentanta $H$. Z twierdzenia Fubiniego dla pewnego
$t_0\in\R$ zachodzi
$H(\cdot,t_0)=f(\cdot)\cosh(\sqrt{\lambda_1}t_0)$
prawie wszędzie. Dzieląc przez dodatni czynnik,
otrzymujemy gładkiego reprezentanta $f$.
Ciągłość wymusza, że jego równanie dystrybucyjne
jest równaniem punktowym.

Dla $f=w$ reprezentant $H$ jest nieujemny wszędzie
(z ciągłości i nieujemności prawie wszędzie).
Gdyby $w$ miała zero w $\Omega$, to ciągłe funkcje
$H(x,t)$ i $w(x)\cosh(\sqrt{\lambda_1}t)$, równe
prawie wszędzie, byłyby równe wszędzie i $H$ osiągałaby
minimum zero w spójnym obszarze $\Omega\times\R$.
Mocna zasada maksimum dla funkcji harmonicznych
~\ref{thm:pde14-strong} dawałaby $H=0$, sprzecznie
z $u\ne0$. Zatem $w>0$ w $\Omega$.
Ponieważ gładki reprezentant $u$ spełnia $|u|=w$
wszędzie, nie ma zer; spójność $\Omega$ wymusza
stały znak $u$.

Przypuśćmy, że przestrzeń własna ma wymiar co najmniej dwa.
Można wtedy znaleźć niezerowe, rzeczywiste, ortogonalne
funkcje własne $u,v$ dla $\lambda_1$.
Po ewentualnej zmianie znaków obie byłyby dodatnie w $\Omega$,
co przeczyłoby $(u,v)_{L^2}=0$: ich iloczyn jest
dodatni na całym obszarze i całkowalny z nierówności
Cauchy'ego--Schwarza. Pierwsza wartość własna jest
więc prosta w rzeczywistej przestrzeni Hilberta,
a po zespoleniu także w zespolonej, gdyż części
rzeczywista i urojona funkcji własnej są rozwiązaniami
tego samego rzeczywistego równania.
\end{proof}

\section{Warunek Neumanna i problem zerowej energii}
\label{sec:pde18-neumann}
\begin{twierdzenie}[Widmo Neumanna i pierwsza dodatnia wartość]
\label{thm:pde18-neumann}
Niech $\Omega\subset\R^d$ będzie ograniczonym, spójnym
obszarem Lipschitzowskim, a $-\Delta_N$ laplasjanem
związanym z formą $\int_\Omega|\nabla u|^2\,
\mathrm d\lambda_d$ o dziedzinie formowej $H^1(\Omega)$.
Jego widmo jest dyskretne i nieujemne:
$0=\mu_1<\mu_2\le\mu_3\le\cdots\to\infty$.
Pierwsza przestrzeń własna składa się z funkcji stałych.
Ponadto
\begin{equation}\label{eq:pde18-neumann}
 \mu_2=\min_{\substack{u\in H^1(\Omega)\setminus\{0\}\\
 \int_\Omega u\,\mathrm d\lambda_d=0}}
 \frac{\int_\Omega|\nabla u|^2\,\mathrm d\lambda_d}
 {\int_\Omega|u|^2\,\mathrm d\lambda_d}>0.
\end{equation}
\end{twierdzenie}
\begin{proof}
Z twierdzenia~\ref{thm:pde10-Rellich-p} zanurzenie
$H^1(\Omega)\hookrightarrow L^2(\Omega)$ jest zwarte.
Forma przesunięta o iloczyn skalarny $L^2$ jest
dokładnie iloczynem skalarnym przestrzeni $H^1$.
Przestrzeń ta jest gęsta w $L^2$, ponieważ zawiera
$C_c^\infty(\Omega)$. Wszystkie kroki konstrukcji
z dowodu twierdzenia~\ref{thm:pde18-compact}
działają z $H^1$ zamiast $H_0^1$ oraz przesunięciem
$c=1$. Otrzymujemy operator samosprzężony o zwartej
rezolwencie, bazę własną oraz zasadę min--max
z dowodu twierdzenia~\ref{thm:pde18-minmax}.

Funkcje stałe należą do $H^1(\Omega)$ i spełniają
równanie formowe z prawą stroną zero. Odwrotnie,
jeżeli $u$ jest funkcją własną dla zera, testowanie
funkcją $u$ daje $\|\nabla u\|_2^2=0$.
Wniosek~\ref{cor:pde10-Poincare} daje
$\|u-u_\Omega\|_2\le C_\Omega\|\nabla u\|_2=0$,
gdzie $u_\Omega=|\Omega|^{-1}\int_\Omega u\,\mathrm d\lambda_d$.
Zatem zero jest prostą wartością własną.
Wszystkie wartości własne są nieujemne, ponieważ
energia jest nieujemna. W drugiej postaci zasady min--max
dla $k=2$ wybieramy osiągającą maksimum przestrzeń
$F=\operatorname{span}\{1\}$, czyli pierwszą przestrzeń
własną. Ortogonalność do niej jest równoważna zerowej
średniej. Minimum osiąga druga funkcja własna;
nierówność Poincarégo--Wirtingera daje dolną granicę
$C_\Omega^{-2}>0$ dla jej ilorazu Rayleigha.
Otrzymujemy~\eqref{eq:pde18-neumann}.
\end{proof}

\begin{uwaga}[Pułapka z monotonicznością]
Monotoniczność wartości własnych względem zawierania
obszarów udowodniliśmy dla warunku Dirichleta.
Nie wolno bez dodatkowych hipotez przenosić jej na
laplasjan Neumanna: przedłużenie przez zero nie działa
wtedy jako zanurzenie $H^1(\Omega_1)\to H^1(\Omega_2)$.
\end{uwaga}

\section{Ślad półgrupy cieplnej i odczytywanie wymiaru}
\label{sec:pde18-heat}
\begin{twierdzenie}[Wiodący człon śladu cieplnego]
\label{thm:pde18-heat}
Niech $\Omega$ spełnia założenia
Twierdzenia~\ref{thm:pde18-weyl}. Dla każdego $t>0$
operator $e^{t\Delta_D}$ jest klasy śladowej w $L^2(\Omega)$,
a przy $t\downarrow0$ zachodzi
\begin{equation}\label{eq:pde18-heat}
 \operatorname{Tr}(e^{t\Delta_D})=
 \sum_{k=1}^\infty e^{-t\lambda_k}
 \sim\frac{|\Omega|}{(4\pi t)^{d/2}}.
\end{equation}
\end{twierdzenie}
\begin{proof}
Prawo Weyla daje $N(\Lambda)\le C(1+\Lambda^{d/2})$.
Dla $t>0$ suma $\sum_ke^{-t\lambda_k}$ jest więc skończona.
Ponieważ $N(\Lambda)=\sum_k\chi_{[\lambda_k,\infty)}(\Lambda)$,
zasada Tonellego prowadzi do tożsamości
\begin{equation}\label{eq:pde18-heat-count}
 \sum_ke^{-t\lambda_k}=t\int_0^\infty e^{-t\Lambda}
 N(\Lambda)\,\mathrm d\Lambda
 =\int_0^\infty e^{-s}N(s/t)\,\mathrm ds.
\end{equation}
Oznaczmy $C_\Omega=\omega_d|\Omega|/(2\pi)^d$.
Z~\eqref{eq:pde18-weyl} wynika, że dla ustalonego $s>0$
$t^{d/2}N(s/t)\to C_\Omega s^{d/2}$.
Wskutek oszacowania liczby modów funkcje pod całką są
zdominowane przez $C e^{-s}(1+s^{d/2})$ dla $0<t\le1$.
Stosujemy zbieżność zdominowaną:
\[
 \lim_{t\downarrow0}t^{d/2}\sum_ke^{-t\lambda_k}
 =C_\Omega\int_0^\infty e^{-s}s^{d/2}\,\mathrm ds
 =C_\Omega\Gamma(1+d/2).
\]
Równość $\omega_d=\pi^{d/2}/\Gamma(1+d/2)$ daje
$C_\Omega\Gamma(1+d/2)=|\Omega|/(4\pi)^{d/2}$.
To dowodzi wzoru~\eqref{eq:pde18-heat}.
\end{proof}

\begin{uwaga}[Powiązanie z geometrią spektralną]
Z zachowania $\operatorname{Tr}(e^{t\Delta_D})\asymp t^{-d/2}$
można odczytać wymiar $d$, a z wiodącego współczynnika
objętość obszaru. Wyższe człony rozwinięcia dla gładkiego
brzegu zawierają jego pole i dane krzywiznowe, ale ich
wyznaczenie wymaga osobnych oszacowań jądra ciepła.
Nie wnioskujemy tu automatycznie z prawa Weyla o pełnym
rozwinięciu asymptotycznym w potęgach $t^{1/2}$.
\end{uwaga}


\begin{lemat}[Przejście abelowe od liczby modów do śladu]
\label{lem:pde18-abel}
Niech $(\nu_j)_{j\ge1}$ będzie niemalejącym ciągiem
liczb dodatnich, $\nu_j\to\infty$, oraz
$N(\Lambda)=\#\{j:\nu_j\le\Lambda\}$. Załóżmy, że dla
pewnych $a>0$, $C>0$, $0\le\gamma<a$ i $K>0$ zachodzi
$|N(\Lambda)-C\Lambda^a|\le K(1+\Lambda^\gamma)$
dla wszystkich $\Lambda\ge0$. Wtedy, gdy $t\downarrow0$,
\[
 \sum_{j\ge1}e^{-t\nu_j}
 =C\Gamma(a+1)t^{-a}+O(1+t^{-\gamma}).
\]
\end{lemat}
\begin{proof}
Ponieważ $N(\Lambda)\le C'(1+\Lambda^a)$, szereg
po lewej jest zbieżny dla $t>0$. Z twierdzenia Tonellego
oraz $e^{-t\nu_j}=t\int_{\nu_j}^\infty e^{-t\Lambda}
\,\mathrm d\Lambda$ otrzymujemy
$\sum_j e^{-t\nu_j}=t\int_0^\infty e^{-t\Lambda}
N(\Lambda)\,\mathrm d\Lambda$.
Po podstawieniu $s=t\Lambda$ wyraz główny ma wartość
$C t^{-a}\int_0^\infty e^{-s}s^a\,\mathrm ds
=C\Gamma(a+1)t^{-a}$. Moduł reszty nie przekracza
$K\int_0^\infty e^{-s}(1+t^{-\gamma}s^\gamma)
\,\mathrm ds=K(1+\Gamma(\gamma+1)t^{-\gamma})$.
\end{proof}

\begin{uwaga}[Różnica między wiodącym wyrazem a rozwinięciem]
Prawo Weyla~\eqref{eq:pde18-weyl} daje tylko $o(t^{-d/2})$
w~\eqref{eq:pde18-heat}. Oszacowanie liczby modów z jawną
resztą prowadzi, przez poprzedni lemat, do oszacowania
reszty śladu. Nie wolno jednak wyprowadzać z samego prawa
Weyla ani dwuczłonowej asymptotyki śladu, ani współczynnika
zależnego od pola brzegu.
\end{uwaga}

\section{Perturbacje widma i wzór Hellmanna--Feynmana}
\label{sec:pde18-perturb}
\begin{wniosek}[Stabilność względem potencjału]
Niech $V_1,V_2\in L^\infty(\Omega;\R)$, a macierz $A$
niech będzie ustalona, symetryczna i jednostajnie eliptyczna.
Wtedy
$|\lambda_k(L_{A,V_1})-\lambda_k(L_{A,V_2})|
\le\|V_1-V_2\|_{L^\infty(\Omega)}$ dla każdego $k\ge1$.
\end{wniosek}
\begin{proof}
Dla każdego $u$ zachodzi
$|q_{V_1}[u]-q_{V_2}[u]|
\le\|V_1-V_2\|_\infty\|u\|_2^2$.
Stosujemy dwukrotnie zasadę min--max.
\end{proof}

\begin{stwierdzenie}[Wzór Hellmanna--Feynmana przy różniczkowalnej gałęzi]
\label{prop:pde18-hf}
Niech $L_t=L_0+tW$, gdzie $L_0$ jest samosprzężonym
operatorem z dyskretnym widmem, a $W$ ograniczonym,
samosprzężonym operatorem na $L^2(\Omega)$.
Przyjmijmy, że na otwartym przedziale $I\subset\R$
można wybrać różniczkowalne w normie grafowej funkcje
$t\mapsto\lambda(t)\in\R$ i
$t\mapsto\varphi(t)\in D(L_0)$ takie, że
$L_t\varphi(t)=\lambda(t)\varphi(t)$ oraz
$\|\varphi(t)\|_2=1$. Wtedy
\begin{equation}\label{eq:pde18-hf}
 \lambda'(t)=(W\varphi(t),\varphi(t))_{L^2}.
\end{equation}
Jeżeli $W$ jest operatorem mnożenia przez funkcję $W(x)$,
to prawa strona ma postać
\[
\int_\Omega W(x)|\varphi(t,x)|^2\,\mathrm d\lambda_d(x).
\]
\end{stwierdzenie}
\begin{proof}
Różniczkujemy równość
$L_t\varphi(t)=\lambda(t)\varphi(t)$ w $L^2$.
Bierzemy iloczyn skalarny z $\varphi(t)$.
Samosprzężoność $L_t$ powoduje skasowanie składników
z $\varphi'(t)$, natomiast z warunku normalizacji
wynika $\|\varphi(t)\|_2=1$. Pozostaje~\eqref{eq:pde18-hf}.
\end{proof}

\begin{uwaga}[Dlaczego założenie różniczkowalności jest jawne?]
Wzór nie mówi sam przez się, że dowolna numerowana gałąź
$\lambda_k(t)$ jest różniczkowalna w punkcie zderzenia
wartości własnych. Dla prostych, izolowanych wartości własnych
regularność gałęzi wynika z teorii perturbacji, której pełnego
dowodu nie wprowadzamy w tej części skryptu.
\end{uwaga}

\section{Podsumowanie i mapa dalszych związków}
\label{sec:pde18-summary}
Z \eqref{eq:pde18-minmax}, \eqref{eq:pde18-weyl} i
\eqref{eq:pde18-liyau} otrzymujemy trzy wzajemnie
uzupełniające się opisy: porządek geometryczny wartości
własnych, asymptotyczną gęstość wysokich częstotliwości
i globalne, ilościowe dolne ograniczenie sumy wartości
własnych. Wzór~\eqref{eq:pde18-heat} pokazuje, że to samo
widmo kontroluje krótkoczasowe zachowanie półgrupy cieplnej.
W następnym rozdziale przejdziemy do pojemności
Sobolewa, osobliwości usuwalnych i własności subtelnych
zbiorów brzegowych.
% Rozdział 19: włączany przed \appendix do PDE_tom_I_wyklad.tex.
